\documentclass[11pt]{article}
\usepackage[margin=1in]{geometry}
\usepackage{amsmath, amssymb, amsthm}
\usepackage{algorithm}
\usepackage{algpseudocode}
\usepackage{booktabs}

\title{Problem 462: Permutation of 3-smooth Numbers}
\author{}
\date{}

\newtheorem{theorem}{Theorem}
\newtheorem{lemma}[theorem]{Lemma}

\begin{document}
\maketitle

\section{Problem Statement}

A \emph{3-smooth number} is an integer whose only prime factors are 2 and 3. For an integer $N$, let $S(N)$ be the set of 3-smooth numbers $\le N$. Define $F(N)$ as the number of permutations of $S(N)$ in which each element appears after all of its proper divisors.

Given: $F(6) = 5$, $F(8) = 9$, $F(20) = 450$, $F(1000) \approx 8.8521816557 \times 10^{21}$.

Find $F(10^{18})$ in scientific notation rounded to ten decimal digits.

\section{Mathematical Foundation}

\begin{theorem}[Poset isomorphism]
\label{thm:poset}
The set of 3-smooth numbers under divisibility is isomorphic as a poset to $\mathbb{Z}_{\ge 0}^2$ under componentwise~$\le$, via $2^a \cdot 3^b \mapsto (a, b)$.
\end{theorem}

\begin{proof}
The map is bijective. For divisibility: $2^{a_1} 3^{b_1} \mid 2^{a_2} 3^{b_2}$ iff $a_1 \le a_2$ and $b_1 \le b_2$, which is precisely the componentwise partial order.
\end{proof}

\begin{lemma}[Young diagram structure]
\label{lem:young}
The set $\{(a, b) : 2^a \cdot 3^b \le N\}$ forms a Young diagram $\lambda = (\lambda_0, \lambda_1, \ldots, \lambda_m)$ with $\lambda_i = \lfloor \log_2(N/3^i) \rfloor + 1$ and $m = \lfloor \log_3 N \rfloor$. The row lengths are non-increasing.
\end{lemma}

\begin{proof}
Row $i$ (corresponding to $b = i$) has cells for $a = 0, 1, \ldots, \lfloor \log_2(N/3^i) \rfloor$, giving $\lambda_i = \lfloor \log_2(N/3^i) \rfloor + 1$ cells. Since $N/3^{i+1} < N/3^i$, we have $\lambda_{i+1} \le \lambda_i$, so the row lengths are weakly decreasing and form a valid Young diagram.
\end{proof}

\begin{theorem}[Frame--Robinson--Thrall hook-length formula]
\label{thm:hlf}
The number of standard Young tableaux of shape $\lambda$ with $n = |\lambda|$ cells is
\[
f^\lambda = \frac{n!}{\displaystyle\prod_{(i,j) \in \lambda} h(i,j)}
\]
where $h(i,j) = (\lambda_i - j) + (\lambda'_j - i) - 1$ is the hook length at cell $(i,j)$, and $\lambda'$ is the conjugate partition.
\end{theorem}

\begin{proof}
Classical result of Frame, Robinson, and Thrall (1954). The proof proceeds by induction on $n$: removing any corner cell reduces to a smaller diagram, and the hook-length product satisfies the same branching recurrence as the count of standard tableaux. See \emph{Canadian J.\ Math.}\ 6 (1954), 316--324.
\end{proof}

\begin{lemma}[Logarithmic computation]
\label{lem:log}
For large $n$, $F(N)$ is computed as
\[
\log_{10} F(N) = \sum_{k=1}^{n} \log_{10} k \;-\; \sum_{(i,j) \in \lambda} \log_{10} h(i,j).
\]
Extracting the integer and fractional parts gives the scientific notation.
\end{lemma}

\begin{proof}
Immediate from taking $\log_{10}$ of the hook-length formula. Writing $\log_{10} F = e + m$ with $e \in \mathbb{Z}$ and $0 \le m < 1$ gives $F = 10^m \times 10^e$.
\end{proof}

\section{Algorithm}

\begin{algorithm}[H]
\caption{Compute $F(N)$ via hook-length formula}
\begin{algorithmic}[1]
\State $m \gets \lfloor \log_3 N \rfloor$
\For{$i = 0$ to $m$}
    \State $\lambda_i \gets \lfloor \log_2(N/3^i) \rfloor + 1$
\EndFor
\State $n \gets \sum_{i} \lambda_i$
\State Compute conjugate $\lambda'_j = |\{i : \lambda_i > j\}|$
\State $\log F \gets \sum_{k=1}^{n} \log_{10} k - \sum_{(i,j) \in \lambda} \log_{10} h(i,j)$
\State Extract mantissa and exponent from $\log F$
\end{algorithmic}
\end{algorithm}

\section{Complexity Analysis}

\begin{itemize}
    \item \textbf{Time:} $O(n)$ where $n = |S(N)| \sim \frac{(\ln N)^2}{2 \ln 2 \cdot \ln 3}$. For $N = 10^{18}$, $n \approx 1107$.
    \item \textbf{Space:} $O(n)$ for storing $\lambda$ and $\lambda'$.
\end{itemize}

\section{Verification}

\begin{center}
\begin{tabular}{crl}
\toprule
$N$ & $|S(N)|$ & $F(N)$ \\
\midrule
6 & 5 & 5 \\
8 & 6 & 9 \\
20 & 10 & 450 \\
1000 & 30 & $\approx 8.8521816557 \times 10^{21}$ \\
\bottomrule
\end{tabular}
\end{center}

\section{Answer}

\[
\boxed{5.5350769703e1512}
\]

\end{document}
